\section{Discussion}

\subsection{Where the Value Lies and Why}

In this simulation, the economic value of a smartphone phenotyping tool depends on whether the tool shortens the breeding cycle. Cycle compression is an assumed scenario input: we found no study that measured how far a phenotyping tool shortens a breeding cycle, so every value here is conditional on an assumed compression of zero to three years. Within that range, cycle compression has a total Sobol index of \NumSobolSTCyc{} for net present value (NPV), against \NumSobolSTErrRed{} for measurement-error reduction. Learning the true compression is worth \NumEVPPIIndiffCyc{} million USD at the indifference hurdle, against \NumEVPPIIndiffErrRed{} million USD for learning the error reduction (Figure~\ref{fig:fig5_what_to_measure}). Across the explored tool profiles, error reduction and cheaper plots raise per-cycle genetic gain by \NumDGdiffPct\% on average (Figure~\ref{fig:fig3_gain_ceiling}), whereas one year of compression raises annual gain by \NumCycOneYearPct\% by construction. At the central plot-level heritability of \NumCfgHtwoYield{}, the per-cycle uplift of a profile with the evidence-based error reduction equals about \NumExchMoHtwoMidErrTen months of cycle compression, and the uplift of a profile with the largest error reduction explored equals about \NumExchMoHtwoMidErrFifty months (Figure~\ref{fig:fig2_time_beats_accuracy}b). These equivalents are computed from the total per-cycle uplift of the robustness block, which also contains the effect of cheaper plots, so they overstate the share due to accuracy alone.

The ordering follows from the breeder's equation and from the model structure, so part of the dominance is structural rather than empirical. Annual gain is per-cycle gain divided by cycle length. Error reduction enters through the square root of heritability, and replication across plots and environments already averages much of the plot error \cite{Cobb2019Enhancing}. Cheaper plots convert into more entries and environments under a fixed budget, with diminishing returns \cite{Cobb2019Enhancing,Lorenz2013Resource}. Compression, in contrast, is a pure time shift in the model. It leaves stages, trial sizes and selection decisions unchanged, so the shorter cycle carries no penalty in selection accuracy or genetic variance. In the economic layer, each program's gain from one cycle accrues over its cycle and then holds constant, with later cycles not credited, so compression enters as an earlier release, a steeper accrual and an earlier adoption path inside the \NumCfgHorizon-year horizon. The horizon-dependent Sobol index of compression shows this timing effect (Figure~\ref{fig:fig4_value_follows_time}c). It is \NumDynSTCycYearFive{} before any release, when only cost phasing differs, peaks at \NumDynSTCycPeak after the first releases of the compressed program, and falls to \NumDynSTCycFinal{} at the horizon as price, discount rate and adoption ceiling gain weight.

The robustness analyses test the structure and the prior, not the existence of compression. The ranking survived a compression of at most one year or at half-year steps, a prior with probability \NumCfgZeroMassB{} on no compression, the removal of the ramp, recurrent cycles, a longer horizon, slower adoption and other calibration targets (Table~\ref{tab:robust}). The least favorable combination still gave a total index of \NumRvSTCycConservativeZeroOne{} for compression against \NumRvSTErrConservativeZeroOne{} for error reduction. These checks show that the ranking does not depend on the width of the assumed range or on the single-cycle ramp. They do not show that a tool can remove a season from the pipeline, and a compression of a few months would make accuracy and compression comparable, as the month equivalents above indicate.

\subsection{Agreement with Earlier Evidence}

The ordering agrees with breeding-program economics. Shorter cycles raised gain per year or lowered the cost of a unit of gain in analyses and simulations of rice, maize and dry bean programs \cite{Bonnecarrere2019Economic,Chaingeni2025More,Lin2023Simulations,Seck2024Stochastic}. Those studies evaluated breeding methods such as rapid generation advance, doubled haploids, speed breeding and rapid recycling. Our result transfers their logic to a phenotyping tool and holds only if the tool shortens the cycle in a comparable way. The small accuracy effect has an empirical counterpart in maize. Image-based tassel traits had lower heritability than manual traits, yet no loss of association-mapping power was detectable \cite{Gage2018Comparing}, which suggests that a measurement with lower heritability does not always degrade genetic inference. The comparison is limited, because the image traits came mostly from one environment and the manual traits from three, so their heritabilities are not directly comparable. The only economic model of phenotyping adoption that we found is theoretical, and its authors reported scarce data on cost and on the contribution of the tool to genetic gain \cite{Awada2018Adoption}. The simulation here supplies a gain contribution, conditional on its inputs.

Two caveats from the literature bound the result. Shorter cycles deplete genetic variance and selection accuracy faster \cite{Seck2024Stochastic,Jighly2019Boosting}, and under a fixed budget more cycles per year shrink the population tested in each cycle \cite{Gorjanc2018Optimal}. The model has no recurrent or rapid-cycling scheme with a selection penalty, so these trade-offs lie outside its scope, and its compression values are best read as values before any such penalty.

\subsection{When Accuracy Matters}

The model values accuracy as an effective reduction in the error variance of the selection phenotype for a single yield trait, and a smartphone imaging tool does not measure grain yield. The evidence for error reduction concerns height, counts and disease scores, and the pathway from images to lower yield-plot error is untested. The accuracy results therefore bound the value of an effect that the tool may not have, and the conclusion that accuracy is worth little is more secure than any positive value of accuracy reported here.

Accuracy carries value in three situations: when plot-level heritability is low, for traits that manual scoring measures poorly, and in programs where the tool cannot change the cycle length. At the lowest plot-level heritability explored, the per-cycle uplift of the largest error reduction equals about \NumExchMoHtwoLowErrFifty months, close to one year of compression, whatever the on-station GxE parameter (Figures~\ref{fig:fig2_time_beats_accuracy}b and~\ref{fig:fig3_gain_ceiling}c). No source reports plot-level heritability for this pipeline, so this condition is a scenario that a program can check against its own trial data. Manual scoring also differs between raters. Trained raters typically reach high concordance with a reference, whereas inexperienced raters can fall far lower \cite{Bock2021Plant}, and rater differences can change estimated allele effects, not only noise \cite{Poland2011Beholder}. In the meta-analysis, pooled agreement between tool and manual scores was lowest for disease severity among the trait groups with several estimates (Figure~\ref{fig:fig6_evidence}b). The model has a single yield trait and cannot value accuracy for visually scored traits, so where such traits drive selection, accuracy may carry more value than we report. Finally, when compression is fixed at zero, error reduction has a total Sobol index of \NumSobolCycZeroSTErrRed{}, nearly tied with the discount rate (\NumSobolCycZeroSTDisc) and ahead of price (\NumSobolCycZeroSTPrice). The value apportioned is small, however. The mean equivalent annual value across profiles without compression is \NumConsEAAmean{} USD per hectare per year (standard deviation (SD) \NumConsEAAsd{}).

\subsection{Implications for Tool Developers}

If a tool is to add value, the simulation indicates that it must bring selection decisions forward, and developers can report the quantities that a program-level valuation needs. Time saved per plot does not equal time saved per cycle. Only savings at the stage that limits a decision shorten the cycle, for example data that arrive early enough to recycle parents once first-stage results exist \cite{Seck2024Stochastic}. Handheld field tools save less time than mechanized platforms (Figure~\ref{fig:fig6_evidence}c) and can save none until the workflow is engineered. In one field deployment the first smartphone workflow was slower than visual scoring, and image analysis took back part of the capture-time saving of a handheld wheat system \cite{Lasdun2024Participatory,Walter2019Highthroughput}. Continued phenotyping is also a precondition for rapid cycling, because simulated gain plateaued when the training data were not updated \cite{Pook2025Optimization}. Throughput and cycle length are therefore complements. Tool papers could report the stage and decision the tool serves, the days saved per decision, the whole-plot cost, the error variance against an independent reference rather than the correlation with a manual score alone, and genotype-dependent bias \cite{Zhou2021Identification}. Agreement with a manual score alone cannot support a valuation of a tool's effect on a program.

\subsection{Implications for Pilots and for the Evaluation of Tools}

If cycle compression is achievable, the simulation indicates that a pilot should first measure the compression the tool delivers and measure error reduction second, against an independent reference. The value-of-information ranking supports this order. Learning cycle compression is worth \NumEVPPIIndiffCyc{} million USD at the indifference hurdle, more than any other input, and it remains the most valuable input when an annual running cost makes the decision depend on the information (Table~\ref{tab:netvalue}). Learning the discount rate (\NumEVPPIIndiffDisc{} million USD) or the bean price (\NumEVPPIIndiffPrice{} million USD) is also valuable, but a pilot cannot learn either, so they are context for scenarios. At a hurdle of zero, EVPPI is \NumEVPPIzeroHurdleMax{} million USD for every input, so the go or no-go decision is insensitive to all uncertainty in a model that excludes the cost of the tool. A pilot also cannot observe a full cycle. It should record the dates at which each stage decision becomes possible under the manual and the tool workflows, the stages at which data are the bottleneck, and the time per plot and per decision, and it should project compression from these records. Error reduction cannot be read from a tool-manual correlation. The pooled correlation (\NumMAtwoR) translates into a negative error reduction if the manual reference is treated as error-free, and into a negative one after the manual-scoring coefficients are converted to correlations with the true value, whereas the as-entered coefficients give a positive one. Repeated manual scoring of the same plots is therefore needed to estimate manual reliability. Expected value of perfect partial information is expressed in the units of the calibrated model, and exceeding the cost of research is necessary but not sufficient for research to be worthwhile \cite{Rothery2020Value}. We did not cost a pilot.

The case suggests an evaluation chain for phenotyping tools: tool metrics (agreement, time and cost), then selection gain at the program level under a fixed budget, then gain calibrated to realized progress, then adoption-weighted value, and finally attribution and value of information.

\subsection{Limitations}

The largest limitation is that cycle compression is assumed, not observed, and it enters the model structurally through the denominator of annual gain and through the timing of release. The tool is also assumed to act on selection error, although the simulated trait is yield, which a smartphone tool does not measure, and the evidence for error reduction concerns trait-level measurement error and not yield-plot error. The valuation also treats compression as free of penalty: it leaves selection decisions unchanged, credits one cycle's gain per program and contains no recurrent or rapid-cycling scheme with a selection penalty. These features favor the shorter cycle, and the shallow downside in Figure~\ref{fig:fig4_value_follows_time}d partly follows from them: both programs spend the same budget, tool costs displace plots instead of adding spending, and in the lookup table negative NPV occurs only at zero compression and only where the simulated per-cycle difference is negative. The robustness analyses vary the ramp, the horizon, the adoption rate and the prior on compression (Table~\ref{tab:robust}), and they do not remove the dependence on the assumed compression. The NPV also ignores the cost of developing, hosting and supporting the tool. The scenarios with an annual running cost (Table~\ref{tab:netvalue}) use illustrative levels that bracket a regional share of the order of magnitude reported for a global system, and a program-specific cost estimate is missing.

Second, the model treats error as random plot error that replication averages out. Image error can be genotype-dependent bias that replication does not remove \cite{Zhou2021Identification}, and a single additive yield trait omits multi-trait selection. Third, the simulation covers one cycle with fixed pipeline sizes. The cycle length is an assumption and the model's $L$ is the time from cross to release, whereas the era-trial rates used for calibration reflect many overlapping cycles. Genotype-by-environment interaction enters only through genotype-mean heritability with the parameterization of Equation~(\ref{eq:heff}), and plot-level heritability and the on-station GxE parameter are assumptions explored in the robustness block (Table~\ref{tab:inputs}).

Fourth, calibration uses one common factor (\NumGainScale{}) to match a realized gain of \NumGainTargetPct\% per year estimated from \NumMAoneKclus programs, three of them Brazilian, and the confidence interval of that estimate extends to approximately zero. Regression-based realized gain is itself imprecise \cite{Rutkoski2019Estimation}. The calibration target changes the scale of every monetary value (Table~\ref{tab:robust}). The factor is equivalent to a lower additive genetic standard deviation and keeps the relative uplift of accuracy, whereas a program that gains slowly because of low effective heritability would raise the value of accuracy. The robustness block does not reach such low heritability, because its manual gain stays far above the realized rate. Fifth, the adoption rate and midpoint are fixed assumptions in the main analysis, and the assumed rate implies faster area gain than the roughly one percentage point per year observed for improved beans \cite{Alene2019Measuring}. The simulated program is a regional network whose budget corresponds to about \NumNetProgMin{} to \NumNetProgMax{} national programs at reported maize and rice pipeline costs \cite{Das2025Costing,Katiyar2026Unified}, and its benefits are summed over the same six countries; this scope is an assumption, and a single national program would have a smaller budget and a smaller area. The model has no price or surplus effects and no farmer-level costs. NPV uses calibrated cell means, so Monte Carlo variation within a profile does not enter its distribution. Ex ante valuations tend toward optimism \cite{Hurley2014Reexamining,Flyvbjerg2021Costbenefit}, so the monetary values rank inputs and are not forecasts.

Sixth, the sensitivity analyses are exact for the grid prior, which weights every simulated level equally and is therefore a choice, not a measured distribution of the inputs. Shapley effects with correlated inputs depend on assumed correlations and no longer lie between the first-order and total Sobol indices \cite{Iooss2019Shapley}, PAWN magnitudes depend on its summary statistic \cite{Puy2020Sensitivity}, and the active-subspace check uses a piecewise constant lookup. We therefore compare ranks, not magnitudes, across measures. The four minor technical inputs cannot be ranked, because their indices are below the noise that Layer 1 simulation adds. Seventh, the meta-analyses rest on few studies with high heterogeneity, include preprints in the agreement analysis, report no variances for time savings, pool agreement across traits and platforms, and had incomplete search coverage, so ``not found'' means not found in our searches. The conversion of manual-reliability coefficients assumes parallel raters with independent errors, and the central error reduction comes from the heritability pairs, because route T1 and the converted route T2 give negative medians and route T2 as entered rests on a pool that mixes coefficient types.

\section{Conclusions}

Under an assumed cycle compression of zero to three years, which no study has measured, the simulated economic value of a smartphone phenotyping tool in a common-bean program lies in shortening the breeding cycle and not in measuring more accurately or more cheaply. Cycle compression has a total Sobol index of \NumSobolSTCyc{} for NPV, against \NumSobolSTErrRed{} for error reduction, and it kept the first rank in all \NumRvNscenarios{} robustness scenarios. At the baseline economic state, the median equivalent annual value is \NumEAAcycOne{} USD per hectare per year with one year of compression and \NumEAAcycZero{} USD without it. Learning the achievable compression is worth \NumEVPPIIndiffCyc{} million USD at the indifference hurdle, against \NumEVPPIIndiffErrRed{} million USD for learning the error reduction. These results are conditional on the assumed compression, on error reduction acting on the selection phenotype and not on measured yield, and on a valuation that excludes the cost of the tool. Accuracy mattered in the simulation only at low plot-level heritability, and the model does not address traits that manual scoring measures poorly.

If these conditions hold, the first measurement for a pilot is the cycle compression the tool achieves, and the evaluation chain from tool metrics to selection gain, adoption-weighted value and value of information gives programs and funders a basis beyond agreement with manual scores.
